\subsection{The exact descent of a coefficient basis to component stalks}
\label{subsec:faithful-free-stalk-transport}

We use \(O=\mathcal O_E\), with the specified integral basis
\(b_1,\ldots,b_n\).  The basis condition in this subsection means
that the vectors \(\rho(b_j)x_l\), indexed by the pair \((l,j)\),
are linearly independent and span the whole underlying module.
It is the source's \(\mathrm{HFree}\) condition when
\(\rho=\rho_H\).

\begin{lemma}[Characteristic-zero coefficient basis]
\label{lem:faithful-free-charzero}
Let \(F\) be a field of characteristic zero, let \(W\) be a
finite-dimensional \(F\)-vector space, and let
\(\rho:O\to\operatorname{End}_F(W)\) be a ring homomorphism.
If its multiplication characteristic polynomials are the
\(r\)-th powers of the regular integral multiplication
characteristic polynomials, mapped to \(F[T]\), then \(W\) has
the coefficient basis condition of rank \(r\).
\end{lemma}
\begin{proof}
The source \texttt{CharZeroFreeModule} chooses an integral
primitive element \(a\in O\) for \(E/\mathbb Q\).
To obtain it, multiply a primitive element in \(E\) by a nonzero
integer so that it becomes integral.
For every \(b\in O\), clearing denominators in its rational
polynomial expression in \(a\) gives
\[
                    N b=P(a),\qquad
                    0\ne N\in\mathbb Z,\quad P\in\mathbb Z[T].
 \tag{ST1}\label{eq:faithful-integral-primitive-generator}
\]
The operator \(\rho(a)\) is semisimple.
Indeed the integral minimal polynomial of \(a\), mapped to \(F[T]\),
annihilates it.  The minimal polynomial over \(\mathbb Q\) is
separable, and its mapped polynomial remains separable in
characteristic zero.  An operator annihilated by a squarefree
polynomial is semisimple.  This applies also to the regular
operator on
\[
                       V=(F\otimes_{\mathbb Z}O)^r.
\]
The characteristic-polynomial assumption at \(a\) says that
these two semisimple operators have the same characteristic
polynomial.

Here is the source's comparison of such operators over \(F\),
which does not require \(F\) to be algebraically closed.
Regard each space as an \(F[T]\)-module.
For a simple summand its characteristic polynomial is the monic
irreducible polynomial \(q\) defining that summand, and the
summand is \(F[T]/(q)\): the annihilator of a nonzero generator
is maximal, and its degree equals the vector-space dimension.
Choose a simple summand on one side.  Its irreducible polynomial
divides the common characteristic polynomial; on the other
semisimple side it divides the characteristic polynomial of some
simple summand, hence that summand is another copy of \(F[T]/(q)\).
Split both summands off and cancel their monic polynomial.
Induction on dimension identifies the complements.
This constructs an \(F\)-linear isomorphism \(e:V\to W\)
intertwining multiplication by \(a\).

It then intertwines every integral polynomial in \(a\).
Equation \eqref{eq:faithful-integral-primitive-generator} and the
invertibility of \(N\) in \(F\) imply that it intertwines every
\(b\in O\).  Let \(x_l\) be the image of \(1\) in the \(l\)-th
factor of \(V\).  The images of its regular integral coefficient
basis are precisely \(\rho(b_j)x_l\), proving the assertion.
\end{proof}

\begin{lemma}[Field descent followed by local lifting]
\label{lem:faithful-free-local-descent}
Let \(R\) be a commutative local ring, let \(K\) be a field,
and let \(f:R\to K\) be a local ring homomorphism.
Let \(I\) be a finite index set and
\(\rho:O\to\operatorname{Mat}_I(R)\) a ring homomorphism.
If the base-changed representation on \(K^I\) has the
coefficient basis condition of rank \(r\), then the representation
on \(R^I\) has that same condition.
The map \(f\) is not assumed surjective.
\end{lemma}
\begin{proof}
Put \(F=R/\mathfrak m_R\).
Localness gives \(\ker f=\mathfrak m_R\), so \(f\) factors as
\(R\twoheadrightarrow F\hookrightarrow K\).
We first descend the basis from \(K\) to \(F\), as in
\texttt{FreeModuleNakayamaDescent.int\_descent}.
Write
\[
                  B=F\otimes_{\mathbb Z}O,\qquad
                  S=B\otimes_F K,\qquad M=F^I .
\]
The residue representation makes \(M\) a finite \(B\)-module.
The given coefficient basis on \(K^I\) is exactly an \(S\)-basis
of rank \(r\): expanding coefficients in the tensor base change
of the integral basis shows that independence and spanning over
\(S\) are equivalent to those coefficient conditions over \(K\).
The canonical base-change identifications give
\(S\otimes_B M\simeq K\otimes_F M\simeq K^I\).

The algebra \(S\) is faithfully flat over the finite-dimensional
Artinian \(F\)-algebra \(B\).
Flatness of \(S\otimes_BM\) over \(S\) descends to flatness of
\(M\) over \(B\).
Explicitly, tensor a monomorphism of \(B\)-modules first with
\(M\), then with the faithfully flat \(S\); its resulting map
is injective by flatness of \(S\otimes_BM\), and faithful
flatness detects the original injectivity.
For every prime of \(B\), choose a prime of \(S\) above it.
Base change at these primes identifies the fibre ranks, which
are all \(r\), because the module over \(S\) is free of rank \(r\).

The source now uses its semilocal finite-flat basis theorem.
In this application \(B\) is Artinian: it is the finite product
of its local Artinian factors.  The finite flat module on each
factor is free; its rank is the just computed \(r\).
The associated orthogonal idempotents combine these local bases
to a \(B\)-basis of \(M\) of rank \(r\).
Expanding this basis in the integral tensor basis gives
\(x_{0,1},\ldots,x_{0,r}\in F^I\) satisfying the coefficient
basis condition over \(F\).
This is faithfully flat basis descent, not a choice of
\(F\)-rational coordinates in an arbitrary nonempty open set.

Lift each coordinate of each \(x_{0,l}\) to \(R\), obtaining
\(x_l\in R^I\).
The square matrix whose columns are \(\rho(b_j)x_l\) has
invertible reduction over \(F\).
It is square because the residue coefficient basis has
\(|I|=nr\) elements.
Its determinant is outside \(\mathfrak m_R\), hence is a unit
of \(R\).  The adjugate gives an inverse matrix, proving
both independence and spanning over \(R\).
This is the source's \texttt{lift\_basis} step.
\end{proof}

We now spell out the exact formal-germ transport in which the
preceding lemmas occur.  For a formal cubical germ \(B\) of
dimension \(d\), use the matrix
\[
 M_B(a)=
 \begin{pmatrix}
 R_a&0\\
 -C_a^t&R_{\bar a}^t
 \end{pmatrix}.
 \tag{ST2}\label{eq:faithful-free-block-matrix}
\]
The canonical isomorphism between the \(2d\) coordinate module
and its product presentation intertwines \(M_B(a)\) with
\(\rho_H(a)\).  The action is a ring homomorphism, by the actual
formal identities.  Consequently its characteristic polynomial
is the one in \(\mathrm{HChar}\), and its coefficient basis is
the one in \(\mathrm{HFree}\).

If \(B'\) is the actual coefficient base change of \(B\) along
\(f:R\to R'\), coefficient extraction gives
\[
                         M_{B'}(a)=f(M_B(a)).
 \tag{ST3}\label{eq:faithful-free-base-change-matrix}
\]
Characteristic polynomials therefore commute with base change.
They also descend along an injective \(f\), by injectivity of
\(R[T]\to R'[T]\).
An actual formal coordinate and gauge isomorphism conjugates
the two doubled actions, using the previously proved matrix
\(T\); hence it preserves \(\mathrm{HChar}\).
These are \texttt{hchar\_map}, \texttt{hchar\_of\_map}, and
\texttt{hchar\_iso}.  The following is the exact local
case distinction in \texttt{HodgeFreeEveryStalk.at\_point}.

Let \(t\) be the supplied \(\mathrm{SAB}\) datum on \(X\), with
\(\mathrm{PowFns}(t)\), and let \(z\in X\).
Choose an actual germ \(g\) over the local ring
\(R=\mathcal O_{X,z}\), satisfying \(\mathrm{HChar}(g.B,r)\).
Assume
\[
 \text{for every prime }p\text{ with }p=0\text{ in }\kappa(z),
                         \quad \mathrm{FibPF}(t,z).
 \tag{ST4}\label{eq:faithful-free-stalk-fibpf}
\]
If \(\kappa(z)\) has characteristic zero,
Lemma~\ref{lem:faithful-free-charzero} gives a coefficient basis
for the reduced block representation on \(\kappa(z)^{2d}\);
Lemma~\ref{lem:faithful-free-local-descent} lifts it to \(R\).

If \(\kappa(z)\) has characteristic \(p>0\), use
\(\mathrm{FibPF}(t,z)\) to choose an algebraically closed
representative \(x:\operatorname{Spec}k\to X\) centred at \(z\),
with \(\mathrm{GeomPF}(t,x)\).
The induced \(R\to k\) is local.
It has characteristic \(p\): otherwise \(p\) would be a unit
in \(k\), hence a unit in \(R\) by localness, contradicting its
zero image in \(\kappa(z)\).
The actual germ-base-change producer supplies \(g_k\), and
\eqref{eq:faithful-free-base-change-matrix} transports its
\(\mathrm{HChar}\) identity.
The positive-characteristic geometric producer
\texttt{HeightCharFreeGeometric.hfree\_geom} supplies
\(\mathrm{HFree}(g_k.B,r)\).
Lemma~\ref{lem:faithful-free-local-descent} now descends and lifts
this exact coefficient basis, giving
\(\mathrm{HFree}(g.B,r)\).

Finally let \(\xi:\operatorname{Spec}\Omega\to X\) be the
source's component-generic point with \(\operatorname{char}\Omega=0\),
and assume the actual predicate \(\mathrm{IsCptGen}(\xi)\).
For its source component index set \(\mathrm{Cpt}(\xi)\),
this supplies the injective structural maps
\(\mathcal O_{X,q}\to\Omega\).
Let \(S\) be the supplied actual \(\mathrm{GermSys}(t,\xi,M)\).
Its generic germ \(S.g_0\) and each \(S.g_q\) have their actual
coefficient-base-change and formal-isomorphism comparisons.
The generic \(\mathrm{HChar}(S.g_0.B,r)\) first transfers through
that isomorphism and then descends along
\(\mathcal O_{X,q}\hookrightarrow\Omega\).
Thus every \(S.g_q\) has the same \(\mathrm{HChar}\) exponent \(r\).
Apply the preceding local case distinction with
\eqref{eq:faithful-free-stalk-fibpf} at each such \(q\).
This is exactly
\texttt{GermSystemHeightFree.hfree\_at\_component\_stalks}.

\paragraph{Frozen-source ledger and status.}
The field-descent and lifting producer is
\href{https://github.com/ChristianKudeba/alpoge-rh-verification/blob/alpoge-source-7cbf7a0c/src/Development/WeilBound/CoreStatements/TorusUniformisation/ChowMahlerFamilies/FreeModuleNakayamaDescent.lean}
{\texttt{FreeModuleNakayamaDescent}}
(blob \texttt{988535984bf3c602ecdf1fa614ae8e8c71ccc922}).
Its characteristic-zero companion is
\href{https://github.com/ChristianKudeba/alpoge-rh-verification/blob/alpoge-source-7cbf7a0c/src/Development/WeilBound/CoreStatements/TorusUniformisation/ChowMahlerFamilies/CharZeroFreeModule.lean}
{\texttt{CharZeroFreeModule}}
(\texttt{b1c6e762bc443d5fddb738cdf5f358c4f031cc96}).
The exact matrix, local case distinction, and generic-component
transport are
\href{https://github.com/ChristianKudeba/alpoge-rh-verification/blob/alpoge-source-7cbf7a0c/src/Development/WeilBound/CoreStatements/TorusUniformisation/ChowMahlerFamilies/HodgeFreeEveryStalk.lean}
{\texttt{HodgeFreeEveryStalk}}
(\texttt{64b7c3a2f5bc38ac43b1990ed38d80ea8c66bb6f}),
wrapped by
\href{https://github.com/ChristianKudeba/alpoge-rh-verification/blob/alpoge-source-7cbf7a0c/src/Development/AutomorphicSpectrum/Proofs/GeometryProofs/GermSystemHeightFree.lean}
{\texttt{GermSystemHeightFree}}
(\texttt{56027647fd66f4f01da68e5a6099af999bf3b5f4}).
All four files were checked against the frozen baseline.

The two algebraic lemmas and the matrix transport above are
proved here.  The positive-characteristic local application
uses the separate actual geometric freeness producer; its
quotient-action, gauge, and degree inputs must be joined before
that application is declared fully translated.
Likewise this subsection takes the actual model's
\(\mathrm{PowFns}\), \(\mathrm{FibPF}\), and
\(\mathrm{IsCptGen}\) inputs as supplied by their own model
producers.  It does not infer any of those from semistability,
rank, or a generic-fibre comparison.
